% Part: proof-theory
% Chapter: sequent-calculus
% Section: admissible-derivable

\documentclass[../../../include/open-logic-section]{subfiles}

\begin{document}

\olfileid{pt}{seq}{adm}

\olsection{অনুমোদনযোগ্য ও নিষ্পাদনযোগ্য বিধি}

প্রমাণতত্ত্বের অনেক ফল এই প্রশ্ন নিয়ে,
অথবা এই প্রশ্নের ফল ব্যবহার করে:
কোনো নির্দিষ্ট বিধি দিয়ে যা কিছু
!!{prove}d করা যায়, সেই বিধি ছাড়াও
কি তা !!{prove}d করা যায়? সবচেয়ে
পরিচিত উদাহরণ কর্তন-বর্জন উপপাদ্য।
এতে দেখা যায়, \Cut{}-সহ কোনো
সিকোয়েন্ট ব্যবস্থায় !!{prove}d
করা যায় এমন যেকোনো সিকোয়েন্ট
\Cut ছাড়াও !!{prove}d করা যায়।
এই ধর্ম এত গুরুত্বপূর্ণ যে এর
একটি নাম আছে।

\begin{defn}
সিকোয়েন্ট ব্যবস্থায় বিধি~$R$
\emph{অনুমোদনযোগ্য} যদি সেই ব্যবস্থা
ও $R$ একসঙ্গে ব্যবহার করে প্রমাণযোগ্য
প্রত্যেক সিকোয়েন্ট $R$ ছাড়াও
প্রমাণযোগ্য হয়।
\end{defn}

\begin{ex}\ollabel{ex:land-deriv}
$\LeftR{\land}$ বিধি \Log{G1c} ও
\Log{G3c}-তে ভিন্ন রূপে আছে।
\Log{G1c}-তে দুটি রূপ: একটির
পূর্বধারণায় প্রধান !!{formula}~$!A \land !B$-এর
বাঁ যৌগাংশ~$!A$, অন্যটির
পূর্বধারণায় ডান যৌগাংশ~$!B$:
\[
\Axiom$ !A, \Gamma \fCenter \Delta$
\RightLabel{$\LeftR{\land}_1$}
\UnaryInf$ !A \land !B, \Gamma \fCenter \Delta$
\DisplayProof
\qquad
\Axiom$!B, \Gamma \fCenter \Delta$
\RightLabel{$\LeftR{\land}_2$}
\UnaryInf$!A \land !B, \Gamma \fCenter \Delta$
\DisplayProof
\]
অন্যদিকে \Log{G3c}-তে একটিই বিধি:
\[
\Axiom$ !A, !B, \Gamma \fCenter \Delta$
\RightLabel{$\LeftR{\land}_3$}
\UnaryInf$ !A \land !B, \Gamma \fCenter \Delta$
\DisplayProof
\]
এবার প্রশ্ন: \Log{G3c}-র
$\LeftR{\land}_3$ বিধি কি
\Log{G1c}-তে অনুমোদনযোগ্য?
হ্যাঁ। $\LeftR{\land}_3$-এর
যেকোনো অনুমানকে শুধু \Log{G1c}-র
বিধি দিয়ে সরাসরি অনুকরণ করা যায়।
$\LeftR{\land}_1$ ও
$\LeftR{\land}_2$-এর সঙ্গে
গঠনগত বিধি~$\LeftR{\Contraction}$-ও লাগবে:
\[
\Axiom$!A, !B, \Gamma \fCenter \Delta$
\RightLabel{$\LeftR{\land}_1$}
\UnaryInf$!A \land !B, !B, \Gamma \fCenter \Delta$
\RightLabel{$\LeftR{\land}_2$}
\UnaryInf$!A \land !B, !A \land !B, \Gamma \fCenter \Delta$
\RightLabel{\LeftR{\Contraction}}
\UnaryInf$!A \land !B, \Gamma \fCenter \Delta$
\DisplayProof
\]
% BN-SRC-624 TARGET-CORRECTION-BEGIN
\par\smallskip\noindent\textnormal{\small\textbf{উৎসসংশোধন (BN-SRC-624):} ∧-বাঁ বিধির অনুকরণের শেষ সিকোয়েন্টে উৎসে \( !A\land !B \Gamma \) লেখা; প্রধান সূত্র ও অবশিষ্ট প্রসঙ্গ আলাদা করতে \( !A\land !B,\Gamma \) করা হয়েছে।}
% BN-SRC-624 TARGET-CORRECTION-END
\end{ex}

অন্য বিধি দিয়ে কোনো বিধির অনুমান
অনুকরণ করা তার অনুমোদনযোগ্যতা
দেখানোর একটি উপায়। তবে একমাত্র
উপায় নয়: কিছু বিধি অনুমোদনযোগ্য,
কিন্তু এভাবে অনুকরণযোগ্য নয়।
যেগুলি অনুকরণযোগ্য, সেগুলিকে
\emph{নিষ্পাদনযোগ্য} বলে।

\begin{defn}
বিধি~$R$,
\[
\AxiomC{$S_1$}
\AxiomC{$\dots$}
\AxiomC{$S_n$}
\RightLabel{$R$}
\TrinaryInfC{$S$}
\DisplayProof
\]
ব্যবস্থাটিতে \emph{নিষ্পাদনযোগ্য}
যদি ব্যবস্থার বিধি ও স্বতঃসিদ্ধের
সঙ্গে $S_1$, \dots,~$S_n$ রূপের
অতিরিক্ত প্রারম্ভিক সিকোয়েন্ট
ব্যবহার করে সিদ্ধান্ত-সিকোয়েন্ট~$S$-এর
একটি ছকভিত্তিক !!{proof} থাকে।
\end{defn}

\olref{ex:land-deriv}-এর !!{proof}
দেখায় যে $\LeftR{\land}_3$ হলো
\Log{G1c}-তে নিষ্পাদনযোগ্য বিধি।
এই !!{proof}-টি
$\LeftR{\land}_3$-এর
পূর্বধারণাগুলিকে প্রারম্ভিক
সিকোয়েন্ট ধরে, কেবল \Log{G1c}-র
বিধি দিয়ে, $\LeftR{\land}_3$-এর
সিদ্ধান্তের একটি
ছকভিত্তিক !!{proof}। (\Log{G1c}-র
কোনো স্বতঃসিদ্ধ ব্যবহার হয়নি।)

\begin{prop}\ollabel{prop:lor-G3-adm}
\Log{G3c}-র $\LeftR{\land}$
ও $\RightR{\lor}$ বিধি \Log{G1c}-তে
নিষ্পাদনযোগ্য।
\end{prop}

\begin{proof}
  $\LeftR{\land}$-এর প্রমাণ
  \olref{ex:land-deriv}-এ আছে।
  $\RightR{\lor}$-এর প্রমাণ
  অনুশীলন হিসেবে রইল।
\end{proof}

\begin{prob}
\Log{G3c}-র $\RightR{\lor}$
বিধি \Log{G1c}-তে নিষ্পাদনযোগ্য,
দেখাও।
\end{prob}

\begin{prob}
ধরি, $\liff$ একটি সংজ্ঞায়িত
!!{operator}: $!A \liff !B$ হলো
$(!A \lif !B) \land (!B \lif !A)$-এর
সংক্ষিপ্ত রূপ। দেখাও, এই বিধিগুলি
\begin{enumerate}
\item \Axiom$\Gamma, !A, !B \fCenter \Delta$
\Axiom$\Gamma \fCenter \Delta, !A, !B$
\RightLabel{\LeftR{\liff}}
\BinaryInf$\Gamma, !A \liff !B \fCenter \Delta$
\DisplayProof
\item
\Axiom$!A, \Gamma \fCenter \Delta, !B$
\Axiom$!B, \Gamma \fCenter \Delta, !A$
\RightLabel{\RightR{\liff}}
\BinaryInf$\Gamma \fCenter \Delta, !A \liff !B$
\DisplayProof
\end{enumerate}
(ক)~\Log{G1c} ও (খ)~\Log{G3c}-তে
নিষ্পাদনযোগ্য।
\end{prob}
% BN-SRC-625 TARGET-CORRECTION-BEGIN
\par\smallskip\noindent\textnormal{\small\textbf{উৎসসংশোধন (BN-SRC-625):} প্রথম দ্বিশর্তক বিধি \(\LeftR{\liff}\), অথচ উৎসের সিদ্ধান্তে \( !A\liff !B \) ডান দিকে ছিল। দুটি পূর্বধারণা দ্বিশর্তক সত্য হলে বাম-ব্যবহারের ক্ষেত্র; সিদ্ধান্তে \( \Gamma,!A\liff !B\Sequent\Delta \) করা হয়েছে। দ্বিতীয় \(\RightR{\liff}\) বিধি অপরিবর্তিত।}
% BN-SRC-625 TARGET-CORRECTION-END


\begin{ex}\ollabel{ex:weak-adm}
\Log{G1c}-র $\RightR{\Weakening}$
বিধি,
\[
\Axiom$ \Gamma \fCenter \Delta$
\RightLabel{\RightR{\Weakening}}
\UnaryInf$ \Gamma \fCenter \Delta, !A$
\DisplayProof
\]
\Log{G3c}-তে অনুমোদনযোগ্য।
প্রমাণের ধারণা সহজ: ধরি,
\Log{G3c}-তে পূর্বধারণা
$\Gamma \Sequent \Delta$-র
!!a{proof}~$\pi$ আছে।
$\pi$-র প্রত্যেক সিকোয়েন্টের
ডান দিকে !!{formula}~$!A$ যোগ করি।
স্বতঃসিদ্ধগুলি স্বতঃসিদ্ধই থাকে,
সঠিক অনুমানগুলিও সঠিক থাকে।
নতুন !!{proof}-এর অন্তিম
সিকোয়েন্ট হলো $\Gamma \Sequent
\Delta, !A$।
% BN-SRC-626 TARGET-CORRECTION-BEGIN
\par\smallskip\noindent\textnormal{\small\textbf{উৎসসংশোধন (BN-SRC-626):} উৎসে নতুন প্রমাণের অন্তিম বস্তুটিকে end-formula বলা হয়েছে, কিন্তু প্রদর্শিত \(\Gamma\Sequent\Delta,!A\) একটি সিকোয়েন্ট; অন্তিম সিকোয়েন্ট বলা হয়েছে।}
% BN-SRC-626 TARGET-CORRECTION-END

তবে $\RightR{\Weakening}$
\Log{G3c}-তে নিষ্পাদনযোগ্য নয়।
ধরি উলটোটি সত্য: \Log{G3c}-র
স্বতঃসিদ্ধ ও বিধির সঙ্গে অতিরিক্ত
প্রারম্ভিক সিকোয়েন্ট~$\Gamma \Sequent \Delta$
নিয়ে $\Gamma \Sequent
\Delta, !A$-এর !!a{proof} পাওয়া যায়।
$\RightR{\Weakening}$ নিষ্পাদনযোগ্য
হতে হলে এই !!{proof} সম্পূর্ণ
ছকভিত্তিক হতে হবে। তখন $\Gamma$
ও~$\Delta$ বাদ দিয়ে সেটি
$\Sequent !A$-এর !!a{proof}-এ
পরিণত করা যেত। অতিরিক্ত প্রারম্ভিক
সিকোয়েন্ট~$\Gamma \Sequent \Delta$
তখন শূন্য সিকোয়েন্ট~$\quad \Sequent \quad$
হত। শূন্য সিকোয়েন্টে কোনো
!!{formula}s নেই, অথচ
\Log{G3c}-র প্রতিটি বিধির
পূর্বধারণায় অন্তত একটি সক্রিয়
!!{formula} লাগে। তাই এই
ছকভিত্তিক !!{proof}-এ শূন্য
সিকোয়েন্ট কোনো অনুমানের
পূর্বধারণা হতে পারে না। ফলে
অতিরিক্ত $\Gamma \Sequent \Delta$
প্রারম্ভিক সিকোয়েন্টটি ব্যবহার
না করলেই কেবল !!{proof}-টি ছকভিত্তিক
হতে পারে। অন্য কথায়, কেবল
\Log{G3c}-র স্বতঃসিদ্ধ ও বিধিতে
$\Gamma \Sequent \Delta, !A$-এর
ছকভিত্তিক নিষ্পাদন থাকত; এই রূপের
প্রত্যেক সিকোয়েন্ট \Log{G3c}-তে
!!a{proof} পেত। কিন্তু
\Log{G3c} বিশুদ্ধ:
সব !!{structure}-তে সত্য
সিকোয়েন্টই কেবল !!{prove}
করতে পারে। $\Gamma = \Delta
= \emptyset$ এবং
$!A \ident \lfalse$ নিলে
\Log{G3c}-কে যেমন
$\quad \Sequent \lfalse$
!!{prove} করতে হত; বাস্তবে
তা পারে না।
\end{ex}

\Log{G3c}-তে শিথিলীকরণ যে
অনুমোদনযোগ্য, তার একটি আরোহী
প্রমাণ এবার দিই। স্বজ্ঞাত যুক্তিটি
আসলে সামান্য শক্তিশালী ফল দেখায়:
$\pi$-র প্রত্যেক সিকোয়েন্টের
ডান দিকে $!A$ যোগ করলেও
!!{proof}-এর গঠন বদলায় না,
বিশেষত উচ্চতা বাড়ে না। এটি
গুরুত্বপূর্ণ বলে আলাদা সংজ্ঞা দিই।

\begin{defn}
বিধি~$R$,
\[
\AxiomC{$S_1$}
\AxiomC{$\dots$}
\AxiomC{$S_n$}
\RightLabel{$R$}
\TrinaryInfC{$S$}
\DisplayProof
\]
ব্যবস্থাটিতে \emph{!!{height}-সংরক্ষণকারী
অনুমোদনযোগ্য} (বা
\emph{!!{hp}-অনুমোদনযোগ্য}) যদি
$i = 1$, \dots,~$n$-এর প্রত্যেকটির
জন্য $\Proves[h_i] S_i$ হলে
$m = \max(h_1, \dots, h_n)$ রেখে
$\Proves_m S$ হয়।
\end{defn}

খেয়াল করো, কোনো বিধি
!!{height}-সংরক্ষণকারী
অনুমোদনযোগ্য হতে এটুকুই যথেষ্ট:
পূর্বধারণা~$S_i$-গুলির
!!{proof}s~$\pi_i$-র
!!{height}~$h_i = \pheight{\pi_i}$
হলে সিদ্ধান্তের এমন
!!a{proof}~$\pi$ থাকবে যার
!!{height}~$\pheight{\pi}$
ওই~$h_i$-গুলির সর্বাধিক মানের
বেশি নয়। বিশেষত একটিমাত্র
পূর্বধারণা~$S_1$ থাকলে
$\pheight{\pi} \le \pheight{\pi_1}$
হলেই চলে। $\RightR{\Weakening}$
ও $\LeftR{\Weakening}$-এর
ক্ষেত্রে আসলে $\pheight{\pi}
= \pheight{\pi_1}$, কিন্তু
সমতা আবশ্যক নয়।

\begin{prop}\ollabel{prop:weak-G3c-adm}
\Log{G1c}-র $\RightR{\Weakening}$
ও $\LeftR{\Weakening}$ বিধি
\Log{G3c}-তে !!{height}-সংরক্ষণকারী
অনুমোদনযোগ্য।
\end{prop}

\begin{proof}
  দেখাতে হবে, $\Log{G3c} \Proves
  \Gamma \Sequent \Delta$-এর
  !!{proof}~$\pi$-র !!{height}
  $\pheight{\pi} = n$ হলে
  $\Gamma \Sequent \Delta, !A$-এর
  \Log{G3c}-!!{proof}~$\pi'$
  আছে, যার $\pheight{\pi'} = n$।
  $n$-এর উপর আরোহ করি।
  $n=0$ হলে $\pi$ কেবল
  $\Gamma \Sequent \Delta$
  স্বতঃসিদ্ধ—অর্থাৎ কোনো
  আণবিক $!C$ একই সঙ্গে
  $\Gamma$ ও~$\Delta$-তে আছে।
  তাই $\Gamma \Sequent \Delta, !A$-ও
  স্বতঃসিদ্ধ। $n = \pheight{\pi} >0$
  হলে !!{proof}~$\pi$-র
  একটি শেষ অনুমান আছে। শেষ
  বিধি অনুযায়ী ক্ষেত্র ভাগ করি;
  এখানে একটিই উদাহরণ দিই।
  ধরি, শেষ বিধি $\RightR{\land}$।
  তখন $\Delta = \Delta', !B \land !C$;
  শেষ অনুমানের অব্যবহিত
  উপ-!!{proof}s $\pi_1$ ও~$\pi_2$-এর
  অন্তিম সিকোয়েন্ট যথাক্রমে
  $\Gamma \Sequent \Delta', !B$
  ও $\Gamma \Sequent \Delta', !C$।
  অর্থাৎ $\pi$-র রূপ:
  \[
  \AxiomC{}
  \RightLabel{$\pi_1$}
  \Deduce$\Gamma \fCenter \Delta', !B$
  \AxiomC{}
  \RightLabel{$\pi_2$}
  \Deduce$\Gamma \fCenter \Delta', !C$
  \RightLabel{\RightR{\land}}
  \BinaryInf$\Gamma \fCenter \Delta', !B \land !C$
  \DisplayProof
  \]
  আরও, $n =
  \max(\pheight{\pi_1}, \pheight{\pi_2}) + 1$।
  তাই $\pheight{\pi_1} < n$
  ও $\pheight{\pi_2} < n$;
  আরোহের অনুমান প্রযোজ্য।
  $\Gamma \Sequent
  \Delta', !B, !A$ এবং
  $\Gamma \Sequent \Delta', !C, !A$-এর
  !!{proof}s $\pi_1'$ ও~$\pi_2'$
  আছে। $\RightR{\land}$ প্রয়োগে
  $\Gamma \Sequent \Delta',
  !B \land !C, !A$-এর
  !!a{proof}~$\pi'$ পাই:
  \[
  \AxiomC{}
  \RightLabel{$\pi_1'$}
  \Deduce$\Gamma \fCenter \Delta', !B, !A$
  \AxiomC{}
  \RightLabel{$\pi_2'$}
  \Deduce$\Gamma \fCenter \Delta', !C, !A$
  \RightLabel{\RightR{\land}}
  \BinaryInf$\Gamma \fCenter \Delta', !B \land !C, !A$
  \DisplayProof
  \]
  এখানে $\pheight{\pi'} =
  \max(\pheight{\pi_1'}, \pheight{\pi_2'}) + 1 = \max(\pheight{\pi_1},
  \pheight{\pi_2}) + 1 = \pheight{\pi}$।
\end{proof}

\olref{prop:weak-G3c-adm} খুব
কাজে লাগে। তার পুনঃপ্রয়োগের
জন্য একটি লিখনরীতি নিই:
$\pi$ যদি $\Gamma \Sequent \Delta$-র
!!a{proof} হয়, তাহলে বাঁ দিকে
$\Pi$-র সব !!{formula}s যোগে
শিথিলীকরণের অনুমোদনযোগ্যতা
এবং ডান দিকে $\Lambda$-র সব
!!{formula}s যোগে একই ধর্ম
প্রয়োগের ফলকে
$\pi[\Pi \Sequent \Lambda]$
লিখি। এটি $\Gamma, \Pi
\Sequent \Delta, \Lambda$-র
!!a{proof}।

\end{document}
